%%%%%%%%%%  scan idx 363  |  folio 350  %%%%%%%%%%
donc \(\ch(\varphi)\equiv\ch(1+a_n)\ \mathrm{mod.}\ \widehat A^{\ge n+1}\)~; or \(\ch(1+a_n)=\eta(\Log(1+a_n))\)
\[
  =\frac{(-1)^{n-1}}{(n-1)!}\,a_n+\cdots
\]

ii) Par substitution, il suffit de montrer la formule (6.5.3) pour \(A=K[X_i;Y_j]_{i,j\in\mathbb N}\) et pour \(\varphi=1+\sum_{k\ge m}X_k\), \(\psi=1+\sum_{k\ge n}Y_k\). Comme \(K\) est sans torsion, on peut plonger cet anneau dans son tensorisé par \(\mathbb Q\), et on voit ainsi que \(\ch\) est injectif. Posons \(\varphi*\psi=1+Z_p+\cdots\), \(Z_p\) étant le premier terme non nul et \(\neq1\) de \(\varphi*\psi\). La multiplicativité de \(\ch\) nous donne, d'après (6.5.2) :
\[
  1+\frac{(-1)^{p-1}}{(p-1)!}.Z_p+\cdots
  =
  \left(1+\frac{(-1)^{m-1}}{(m-1)!}.X_m+\cdots\right)
  \left(1+\frac{(-1)^{n-1}}{(n-1)!}.Y_n+\cdots\right),
\]
d'où la valeur de \(Z_p\) en identifiant les termes de plus bas degré.

\textbf{6.6.} Le \(\lambda\)-anneau \(\Ch A\) est augmenté vers \(K\)~: on définit en effet un \(\lambda\)-homomorphisme de \(\Ch A\) dans \(K\) par \((a,\varphi)\rightsquigarrow a\). Il est donc muni d'une filtration canonique de \(\lambda\)-anneau augmenté. Nous allons la comparer à celle qui vient d'être définie.

\textbf{Proposition 6.6.1.} Soit \(x\in\Ch A\) un élément de la forme \(x=(n,1+\cdots+a_n)\). Alors :
\[
\begin{array}{ll}
\text{i)}& \lambda^i(x)=0\quad\text{pour }i>n~;\\[0.3em]
\text{ii)}& \lambda^n(x)=(1,1+a_1)\\[0.3em]
\text{iii)}& \displaystyle\sum_{i=0}^n(-1)^i\lambda^i(x)
       =(0,1-(n-1)!a_n+\cdots),
\end{array}
\]
les termes non écrits étant de degré \(>n\).

Par fonctorialité, il suffit de faire la démonstration dans le cas 
